% ===========================================================================
\chapter{Fondements et référentiels}
\label{chap:fondements}
% ===========================================================================

\section{De DevOps à DevSecOps}

DevOps est défini comme une approche organisationnelle qui met l'accent sur la
collaboration transversale au sein et entre les équipes de développement
logiciel, afin d'exploiter des systèmes résilients et d'accélérer la livraison
des changements~\cite{prates2025}. Ses principes sont fréquemment résumés par
l'acronyme CAMS : culture, automatisation, mesure et partage.

DevSecOps étend ces principes en y intégrant la sécurité. Prates et Pereira
montrent que chacun des quatre principes CAMS admet une déclinaison sécurité :
la culture s'étend à la collaboration avec l'équipe sécurité, l'automatisation
englobe l'automatisation des contrôles de sécurité, la mesure intègre des
métriques de menaces et de vulnérabilités, et le partage inclut la mise en
commun des problématiques de sécurité~\cite{prates2025}.

Les auteurs considèrent DevSecOps comme une extension naturelle de DevOps,
position également retenue par la norme IEEE~2675~\cite{ieee2675}.

Le principe structurant est celui du \emph{shift security to left} : inclure les
pratiques de sécurité dès les étapes les plus précoces du cycle de
développement. Ce principe a une conséquence directe sur l'ordonnancement des
travaux, exposée au chapitre~\ref{chap:phase0}.

\section{Le cycle de vie DevSecOps}

Prates et Pereira identifient les phases du flux DevSecOps et y rattachent
chaque pratique~\cite{prates2025}. Les phases sont les suivantes :
\emph{Create}, \emph{Plan}, \emph{Adapt}, \emph{Verify},
\emph{Preproduction}, \emph{Release}, \emph{Predict}, \emph{Respond},
\emph{Detect} et \emph{Configure}.

Cette cartographie présente un intérêt pratique immédiat : elle indique
\emph{où} dans la chaîne chaque pratique doit intervenir. Elle a servi de
squelette à l'architecture de notre pipeline, présentée au
chapitre~\ref{chap:archi-pipeline}.

\section{La taxonomie des pratiques de Prates et Pereira}

\subsection{Méthode et portée de l'étude}

L'étude de référence du module est une revue de littérature multivocale
(\emph{Multivocal Literature Review}) portant sur 103 publications
primaires~\cite{prates2025}. Contrairement à une revue systématique classique
limitée à la littérature académique, une revue multivocale intègre également la
littérature grise : rapports techniques d'entreprises, pages web, livres
blancs.

La répartition des sources est quasiment équilibrée, avec 49\% de littérature
blanche et 51\% de littérature grise, ce que les auteurs interprètent comme une
convergence entre la recherche académique et la pratique industrielle.

L'étude identifie \textbf{39 pratiques DevSecOps distinctes}.

\subsection{Le classement par niveaux}

Les auteurs classent ces 39 pratiques en trois niveaux selon le nombre de
publications les mentionnant~\cite{prates2025} :

\begin{itemize}
  \item \textbf{Tier 1 --- Core} : pratiques mentionnées dans \textbf{30
  publications ou plus}. Elles constituent les pratiques fondamentales de
  DevSecOps.
  \item \textbf{Tier 2 --- Important} : pratiques mentionnées dans
  \textbf{10 à 29} publications. Également considérées comme importantes.
  \item \textbf{Tier 3 --- Nice to have} : pratiques mentionnées dans
  \textbf{moins de 10} publications. Interprétées comme souhaitables mais non
  cruciales pour une mise en œuvre de DevSecOps.
\end{itemize}

Il importe de souligner que le niveau d'une pratique reflète son degré
d'établissement dans la littérature, et \textbf{non} son caractère automatisable
ni sa position dans la pipeline. Cette distinction a des conséquences directes,
développées à la section~\ref{sec:pratiques-non-outillables}.

\subsection{Les onze pratiques de niveau Tier 1}

Le tableau~\ref{tab:tier1} présente les onze pratiques de niveau Tier~1, qui
constituent le socle obligatoire du projet.

\begin{table}[H]
\centering
\caption{Les onze pratiques DevSecOps de niveau Tier~1 (d'après
Prates et Pereira~\cite{prates2025}, tableau~7)}
\label{tab:tier1}
\small
\begin{tabular}{@{}llcl@{}}
\toprule
\textbf{Id.} & \textbf{Pratique} & \textbf{Mentions} & \textbf{Phase du cycle} \\
\midrule
P1  & SAST                        & 65 & Verify \\
P2  & DAST                        & 58 & Verify \\
P3  & Automated security testing  & 56 & Verify \\
P4  & Vulnerability scanning      & 44 & Verify / Preproduction \\
P5  & Continuous monitoring       & 41 & Predict / Detect \\
P6  & Continuous integration      & 41 & Adapt \\
P7  & Threat modelling            & 40 & Plan / Predict \\
P8  & Shift left security         & 36 & Create \\
P9  & Collaboration               & 36 & Create \\
P10 & Continuous deployment       & 34 & Release \\
P11 & Code review                 & 33 & Verify \\
\bottomrule
\end{tabular}
\end{table}

Les définitions retenues, reformulées d'après l'article source, sont les
suivantes.

\begin{description}[leftmargin=1.2cm,style=nextline]
  \item[\pid{P1} SAST] Analyse du code source, du bytecode ou du binaire d'une
  application afin d'identifier des vulnérabilités, sans exécuter le programme.

  \item[\pid{P2} DAST] Test d'une application en cours d'exécution, vue de
  l'extérieur, par envoi de requêtes et analyse des réponses, en simulant des
  attaques réelles.

  \item[\pid{P3} Automated security testing] Recours à des outils et processus
  automatisés pour vérifier systématiquement la sécurité à différentes étapes
  du développement. Couvre les tests d'API, les tests d'applications web et les
  tests unitaires orientés sécurité. Cette pratique est \emph{distincte} de
  SAST et de DAST.

  \item[\pid{P4} Vulnerability scanning] Identification automatisée et continue
  des vulnérabilités du logiciel \emph{et} de l'infrastructure, par
  confrontation à des bases de CVE connues. L'article précise que cette pratique
  englobe l'analyse des dépendances logicielles et celle des images de
  conteneurs.

  \item[\pid{P5} Continuous monitoring] Collecte et analyse automatisées et
  continues des données de performance et de sécurité des systèmes en
  fonctionnement, procurant une visibilité en temps réel.

  \item[\pid{P6} Continuous integration] Intégration automatique et continue
  des modifications de code dans un dépôt partagé, avec construction, test,
  validation et contrôles de sécurité automatisés.

  \item[\pid{P7} Threat modelling] Processus systématique et proactif
  d'identification et d'évaluation des menaces, par analyse de la conception, de
  l'architecture et des fonctionnalités du système.

  \item[\pid{P8} Shift left security] Intégration des considérations de sécurité
  le plus tôt possible dans le cycle de développement, plutôt qu'à la fin.

  \item[\pid{P9} Collaboration] Collaboration transversale entre les équipes de
  développement, d'exploitation et de sécurité, par décloisonnement, afin de
  créer une culture consciente de la sécurité.

  \item[\pid{P10} Continuous deployment] Construction, test et déploiement
  automatiques et continus des modifications vers la production, les contrôles
  de sécurité étant partie intégrante de la chaîne automatisée.

  \item[\pid{P11} Code review] Examen systématique du code par des pairs ou des
  experts afin d'identifier défauts, vulnérabilités et problèmes de qualité
  avant fusion.
\end{description}

\subsection{Pratiques de niveau Tier 2}

Le tableau~\ref{tab:tier2} recense les vingt pratiques de niveau Tier~2 et
indique celles retenues pour ce projet.

\begin{table}[H]
\centering
\caption{Pratiques de niveau Tier~2 et sélection retenue}
\label{tab:tier2}
\small
\begin{tabular}{@{}llcl@{}}
\toprule
\textbf{Id.} & \textbf{Pratique} & \textbf{Mentions} & \textbf{Retenue} \\
\midrule
P12 & Infrastructure as code         & 28 & \fait{} oui \\
P13 & SCA                            & 27 & \fait{} oui \\
P14 & Security training              & 27 & \fait{} oui \\
P15 & VCS                            & 26 & \fait{} oui \\
P16 & Penetration testing            & 24 & \fait{} oui \\
P17 & Continuous compliance          & 22 & \fait{} oui \\
P18 & Logging                        & 21 & \fait{} oui \\
P19 & IAST                           & 18 & \nonfait{} non \\
P20 & Secure coding                  & 18 & \fait{} oui \\
P21 & Configuration management       & 17 & optionnel \\
P22 & Environment hardening          & 16 & \fait{} oui \\
P23 & Risk assessment                & 15 & \fait{} oui \\
P24 & Security requirement gathering & 13 & \fait{} oui \\
P25 & RASP                           & 13 & \nonfait{} non \\
P26 & Security policies              & 12 & optionnel \\
P27 & Patch management               & 12 & \fait{} oui \\
P28 & Security incident management   & 11 & \fait{} oui \\
P29 & IAM                            & 11 & optionnel \\
P30 & Secrets management             & 11 & \fait{} oui \\
P31 & Security auditing              & 11 & \fait{} oui \\
\bottomrule
\end{tabular}
\end{table}

\begin{constat}{trois pratiques fréquemment supposées Tier 1 n'en sont pas}
L'analyse de la source corrige trois idées reçues qu'il importe de ne pas
reproduire dans ce rapport :

\begin{itemize}
  \item \textbf{SCA} (analyse de composition logicielle) est une pratique de
  niveau Tier~2 (\pid{P13}, 27 mentions), et non Tier~1. La pratique de niveau
  Tier~1 qui couvre l'analyse des dépendances et des images de conteneurs est
  \pid{P4} \emph{Vulnerability scanning}.
  \item \textbf{Secrets management} est également de niveau Tier~2
  (\pid{P30}, 11 mentions).
  \item \textbf{Infrastructure as code} est de niveau Tier~2 (\pid{P12},
  28 mentions), à deux mentions du seuil de 30.
\end{itemize}

Cette distinction a une conséquence favorable sur le projet. Le sujet exige par
ailleurs l'emploi d'outils SAST, SCA, \emph{secret scanning} et DAST. SAST et
DAST relevant du Tier~1 et SCA et \emph{secrets management} du Tier~2, leur
mise en œuvre satisfait simultanément une exigence obligatoire du sujet et
ouvre droit à la valorisation des pratiques Tier~2.
\end{constat}

\subsection{Pratiques de niveau Tier 3}

Pour mémoire, les huit pratiques de niveau Tier~3 sont : \emph{Security
champions} (8 mentions), \emph{Red team} (8), \emph{WAF} (8),
\emph{Alerting} (7), \emph{Continuous reporting} (6), \emph{Software
licensing} (3), \emph{Chaos engineering} (3) et \emph{Bug bounties} (2).

Trois d'entre elles sont obtenues à coût marginal quasi nul dans notre
architecture et ont été retenues à ce titre : \emph{Alerting}, fourni par la
pile de supervision ; \emph{WAF}, par un reverse-proxy en frontal ; et
\emph{Software licensing}, par une option de l'outil d'analyse de composition.

\section{Pratiques outillables et pratiques de processus}
\label{sec:pratiques-non-outillables}

Une lecture attentive du tableau~\ref{tab:tier1} fait apparaître un point
méthodologique essentiel : \textbf{quatre des onze pratiques de niveau Tier~1 ne
sont pas outillables} au sens d'un analyseur produisant un rapport.

\begin{table}[H]
\centering
\caption{Nature des pratiques Tier~1 et nature des preuves correspondantes}
\label{tab:nature-preuves}
\small
\begin{tabular}{@{}p{3.4cm}p{4.6cm}p{5.6cm}@{}}
\toprule
\textbf{Nature} & \textbf{Pratiques} & \textbf{Nature de la preuve} \\
\midrule
Automatisable en
pipeline
& \pid{P1} SAST, \pid{P2} DAST,
\pid{P3} tests automatisés,
\pid{P4} scan de vulnérabilités
& Rapport d'outil, journal
d'exécution, échec de build \\
\addlinespace
Automatisable
hors analyse
& \pid{P6} intégration continue,
\pid{P10} déploiement continu
& Définition de la pipeline,
historique d'exécutions \\
\addlinespace
Processus et
organisation
& \pid{P7} modélisation des menaces,
\pid{P8} shift left,
\pid{P9} collaboration,
\pid{P11} revue de code
& Document, configuration de
plate-forme, fil de discussion
de \emph{pull request} \\
\addlinespace
Exploitation
& \pid{P5} supervision continue
& Tableau de bord, alerte
déclenchée \\
\bottomrule
\end{tabular}
\end{table}

Cette répartition explique la formulation du sujet, qui demande d'automatiser
l'exécution des pratiques \emph{autant que possible}. Pour les pratiques de
processus, la preuve attendue n'est pas un rapport d'analyse mais l'artefact
lui-même : le modèle de menaces, les règles de protection de branche, le fichier
d'attribution de propriété du code, ou un fil de revue montrant une
vulnérabilité discutée puis corrigée.

\section{Référentiels de sécurité appliqués}

\subsection{OWASP ASVS 5.0 comme référentiel de vérification}

La taxonomie de Prates et Pereira indique \emph{quelles pratiques} mettre en
œuvre dans la chaîne de production. Elle ne dit rien des propriétés de sécurité
que l'application doit elle-même présenter. Ces deux questions sont distinctes et
appellent deux référentiels distincts.

L'ASVS~5.0~\cite{owaspasvs5} fournit le second. Il définit environ 350 exigences
de sécurité applicative réparties en 17 chapitres, organisées en trois niveaux
de vérification. L'ASVS précise explicitement qu'il ne prescrit ni les activités
du cycle de développement, ni la manière dont l'application doit être construite
par une pipeline CI/CD : il spécifie les résultats de sécurité à atteindre au
sein du produit. Les deux référentiels sont donc complémentaires et non
redondants.

Deux apports de l'ASVS~5.0 ont particulièrement orienté le projet.

Le premier est la notion de \textbf{décision de sécurité documentée}. L'ASVS
introduit des exigences de documentation au motif que certains contrôles
dépendent trop du contexte applicatif pour être prescrits de manière générique,
et que sans décision écrite les exigences de mise en œuvre correspondantes ne
sont tout simplement pas vérifiables.

Le second est la position de l'ASVS sur les limites de l'outillage automatisé.
Le référentiel indique que les outils SAST et DAST, correctement intégrés,
peuvent détecter des problèmes de sécurité, mais qu'ils ne sont pas en mesure de
vérifier un grand nombre des exigences les plus complexes, en particulier celles
relatives à la logique métier et au contrôle d'accès. Il en conclut que des
vérifications spécifiques à l'application doivent être écrites, de nature
comparable aux tests unitaires et d'intégration.

\begin{constat}{convergence des deux référentiels}
Cette position de l'ASVS rejoint exactement la pratique \pid{P3}
\emph{Automated security testing}, troisième pratique la plus citée de l'étude
de Prates et Pereira avec 56 mentions, et définie comme distincte de SAST et de
DAST.

Les deux sources convergent donc vers un même livrable : une suite de tests de
sécurité spécifiques à l'application, portant notamment sur l'autorisation et
la logique métier. Ce livrable fait l'objet du chapitre~\ref{chap:tests}.
\end{constat}

\subsection{Niveau de vérification ASVS retenu}

L'ASVS~5.0 définit trois niveaux. Le niveau~1 représente environ 20\% des
exigences et constitue la première couche de défense. Le niveau~2 correspond au
niveau que la majorité des applications devraient viser. Le niveau~3 s'adresse
aux applications exigeant la plus haute assurance.

Le référentiel encourage explicitement les organisations à produire une version
adaptée à leur domaine, en écartant les chapitres non pertinents, à condition
que le périmètre retenu et les exclusions soient justifiés dans le rapport de
vérification.

Le niveau visé par ce projet et la justification des exclusions sont présentés
au chapitre~\ref{chap:modelisation}.

\subsection{Catalogues de risques}

Les catalogues de risques suivants sont pris en compte, conformément au sujet :

\begin{itemize}
  \item \textbf{OWASP Top 10:2025}~\cite{owasptop10} pour les risques
  applicatifs web ;
  \item \textbf{OWASP API Security Top 10}~\cite{owaspapi} pour les risques
  propres aux API REST, avec une attention particulière à l'authentification, à
  l'autorisation, au contrôle d'accès, à la validation des entrées, à
  l'exposition des données et à la protection contre les abus ;
  \item \textbf{OWASP Top 10 for LLM Applications}~\cite{owaspllm} pour les
  risques liés aux fonctionnalités fondées sur un modèle de langage ;
  \item \textbf{OWASP Mobile Top 10}~\cite{owaspmobile} pour l'application
  mobile, dont le traitement est reporté à la seconde phase du projet.
\end{itemize}
